% =====================================================================
% 第 IV 章  Class invariants II (A. Borel, Nov. 13, 1957)
% =====================================================================
\lecture{Class invariants II}{Class Invariants II}{A. Borel}{Nov. 13, 1957}

\setcounter{theorem}{0}
\setcounter{lemma}{0}

% ---------------------------------------------------------------
\section{Introduction}\label{sec:IV-intro}

The main purpose of this lecture is to obtain Theorem~3. For this we
shall first establish some properties of certain functions formed by
means of the discriminant $\Delta$ (see §§2, 3). We follow the usual
conventions of algebraic number theory in which often no distinction is
made between an ideal $\mathfrak{a}$ in an algebraic number field $K$ and
the ideal $\mathfrak{o}(L)\cdot\mathfrak{a}$ it generates in a finite
extension $L$ of $K$. In particular let $\mathfrak{a}_{1}$,
$\mathfrak{a}_{2}$ be ideals in $K_{1}$, $K_{2}$ and $L$ be a finite
extension of $K_{1}$ and $K_{2}$. Then we say that $\mathfrak{a}_{1}$
divides $\mathfrak{a}_{2}$ (resp.\ that $\mathfrak{a}_{1}$ and
$\mathfrak{a}_{2}$ are prime to each other) if
$\mathfrak{o}(L)\cdot\mathfrak{a}_{1}$ divides
$\mathfrak{o}(L)\cdot\mathfrak{a}_{2}$ (resp.\
$\mathfrak{o}(L)\cdot\mathfrak{a}_{1}$ and
$\mathfrak{o}(L)\cdot\mathfrak{a}_{2}$ are prime to each other). This is
then true in any algebraic extension of $K_{1}$ and $K_{2}$.

% ---------------------------------------------------------------
\section{\texorpdfstring{The functions $\varphi_{M}$}{The functions varphi\_M}}\label{sec:IV-varphiM}

In dealing with $\Delta$, it will be convenient to use the homogeneous
formulation. We denote in the same way the automorphism of the upper half
plane $E$ given by
\[
z\longrightarrow(az+b)(cz+d)^{-1}
\]
and the homogeneous linear transformation
$(w_{1},w_{2})\longrightarrow(aw_{1}+bw_{2},\ cw_{1}+dw_{2})$. We recall
that if $h(z)$ is a modular form of weight $k$, then
\[
h(w_{1},w_{2})=w_{2}^{-2k}h(w_{1}/w_{2})
\]
is a homogeneous function of degree $-2k$ in $w_{1},w_{2}$ and is
invariant under the ``homogeneous'' modular group $G$. In particular, the
modular form $\Delta$ gives rise to a homogeneous function of degree
$-12$, invariant under $G$:
\[
\Delta(w_{1},w_{2})=w_{2}^{-12}\,\Delta\!\left(\frac{w_{1}}{w_{2}}\right),
\]
and we have
\[
\Delta(w_{1},w_{2})=\left(\frac{2\pi}{w_{2}}\right)^{12}
q\,(1+B(q)),\qquad q=e^{2\pi i\frac{w_{1}}{w_{2}}},
\tag{1}\label{eq:IV-1}
\]
where $B(q)$ is a power series of $q$ with integral rational coefficients
and $B(0)=0$ (see II, §6).

For any $M$ in $H_{n}$, put
\[
\varphi_{M}(w_{1},w_{2})=n^{12}\frac{\Delta(M(w_{1},w_{2}))}{\Delta(w_{1},w_{2})}.
\]

Then $\varphi_{M}$ depends only upon the class $GM$ of $H_{n}/G$ and,
from the $N=\psi(n)$ classes of $H_{n}/G$, we obtain $N$ homogeneous
functions of degree $0$ in $w_{1},w_{2}$:
\[
\varphi_{i}(w_{1},w_{2})=\varphi_{M_{i}}(w_{1},w_{2}),\qquad
1\leq i\leq N.
\]

These functions are regular for $\operatorname{Im}\bigl(\frac{w_{1}}{w_{2}}\bigr)>0$
and are permuted among themselves under the transformations in $G$.

Now, it follows immediately from \eqref{eq:IV-1} that
\[
\varphi_{i}(w_{1},w_{2})
=a_{i}^{12}\zeta_{d_{i}}^{b_{i}}q^{\frac{a_{i}}{d_{i}}-1}
\bigl(1+B(\zeta_{d_{i}}^{b_{i}}q)\bigr)\bigl(1+B(q)\bigr)^{-1},
\quad\Bigl(\zeta_{d}=e^{\frac{2\pi i}{d}}\Bigr),
\tag{2}\label{eq:IV-2}
\]
and we see, by a similar argument as in III, §4, that
\[
\Phi_{n}(t,j)=\prod_{i=1}^{N}\bigl(t-\varphi_{i}(w_{1},w_{2})\bigr)
\]
is a polynomial in $t$ and $j$ with integral rational coefficients.

\begin{lemma}\label{lem:IV-1}
If $n=p$ is a prime, $\prod_{i=1}^{N}\varphi_{i}(w_{1},w_{2})$ is a
constant and is equal to $(-1)^{p-1}p^{12}$.
\end{lemma}

\begin{proof}[Proof.]
In this case, $N=\psi(p)=p+1$ and the matrices $M_{i}$ are given by
\[
M_{i}=
\begin{pmatrix}1&i\\ 0&p\end{pmatrix}\quad(1\leq i\leq p);\qquad
M_{p+1}=
\begin{pmatrix}p&0\\ 0&1\end{pmatrix}.
\]
Hence,
\[
\left\{
\begin{aligned}
\varphi_{i}(w_{1},w_{2})
&=\zeta_{p}^{i}\,q^{\frac{1}{p}-1}
\bigl(1+B(\zeta_{p}^{i}q^{\frac{1}{p}})\bigr)
(1+B(q))^{-1},\quad(1\leq i\leq p),\\[3pt]
\varphi_{p+1}(w_{1},w_{2})
&=p^{12}q^{p-1}\bigl(1+B(q^{p})\bigr)(1+B(q))^{-1},
\quad(i=p+1),
\end{aligned}\right.
\tag{3}\label{eq:IV-3}
\]
and the q-expansion of $\prod_{i=1}^{N}\varphi_{i}(w_{1},w_{2})$ starts
with the constant term $(-1)^{p-1}p^{12}$. However, as this product is
also a polynomial in $j$, it must be equal to the constant
$(-1)^{p-1}p^{12}$.
\end{proof}

% ---------------------------------------------------------------
\section{\texorpdfstring{Properties of the singular values of the $\varphi_{M}$}{Properties of the singular values of the varphi M}}\label{sec:IV-singular}

We now fix an imaginary quadratic field $K$.

\begin{lemma}\label{lem:IV-2}
Let $\mathfrak{a}$ be an ideal of $K$ and $(\alpha_{1},\alpha_{2})$ a
basis of $\mathfrak{a}$ such that
$\operatorname{Im}(\alpha_{1}/\alpha_{2})>0$. Then, for any $M$ in
$H_{n}$, $\varphi_{M}(\alpha_{1},\alpha_{2})$ is an algebraic integer
and, if $n=p$ is a prime, the principal ideal
$\bigl(\varphi_{M}(\alpha_{1},\alpha_{2})\bigr)$ is a divisor of
$(p^{12})$.
\end{lemma}

\begin{proof}[Proof.]
By the definition of $\Phi_{n}$,
\[
\Phi_{n}\bigl(t,j(\mathfrak{a})\bigr)
=\prod_{i=1}^{N}\bigl(t-\varphi_{i}(\alpha_{1},\alpha_{2})\bigr).
\]
But, as $j(\mathfrak{a})$ is an algebraic integer (III, Theorem 2), the
left hand side is a unitary polynomial in $t$ with integral algebraic
coefficients. Since
$\varphi_{M}(\alpha_{1},\alpha_{2})$
($=\varphi_{i}(\alpha_{1},\alpha_{2})$ for some $i$) is a root of
$\Phi_{n}(t,j(\mathfrak{a}))=0$, it is an algebraic integer. The second
part is then an immediate consequence of Lemma~\ref{lem:IV-1}.
\end{proof}

Let $\mathfrak{a}$ be an integral ideal of $K$. Then, for any (fractional
or integral) ideal $\mathfrak{b}$ of $K$, $\mathfrak{a}\cdot\mathfrak{b}$
is an additive subgroup of $\mathfrak{b}$, of index equal to $N\mathfrak{a}$,
as follows from the fact that
$N(\mathfrak{a}\cdot\mathfrak{b})=N(\mathfrak{a})\,N(\mathfrak{b})$. Thus,
if $(\beta_{1},\beta_{2})$ is a base of $\mathfrak{b}$, and
$(a\beta_{1}+b\beta_{2},\ c\beta_{1}+d\beta_{2})$ a base of
$\mathfrak{a}\cdot\mathfrak{b}$, where $a,b,c,d\in\mathbb{Z}$, we must
have $ad-bc=N\mathfrak{a}$; moreover $(a,b,c,d)=1$ if $\mathfrak{a}$ is
not divisible by a rational integer $>1$, for instance if $\mathfrak{a}$
is prime.

A prime ideal $\mathfrak{p}$ of $K$ is said to be of first degree if
every integer of $K$ is congruent mod $\mathfrak{p}$ to a rational number,
or equivalently if $N\mathfrak{p}=p$ is a prime number; $\mathfrak{p}$ is
unramified for $K/\mathbb{Q}$ if its square does not divide a prime
rational number, or equivalently if $N\mathfrak{p}$ does not divide the
discriminant of $K/\mathbb{Q}$. Under those two conditions we have
$\mathfrak{p}\neq\bar{\mathfrak{p}}$ ($\bar{\phantom{x}}$ denoting complex
conjugation) and $(p)=\mathfrak{p}\cdot\bar{\mathfrak{p}}$.

\begin{theorem}\label{thm:IV-1}
Let $(\alpha_{1},\alpha_{2})$ be a basis of an ideal $\mathfrak{a}$ of
$K$, $\mathfrak{p}$ be a prime ideal of the first degree in $K$,
unramified for $K/\mathbb{Q}$, and $p=N\mathfrak{p}$. Let $P\in H_{p}$
(resp.\ $\bar{P}\in H_{p}$) be such that $P(\alpha_{1},\alpha_{2})$
(resp.\ $\bar{P}(\alpha_{1},\alpha_{2})$) is a basis of
$\mathfrak{p}\cdot\mathfrak{a}$ (resp.\ $\bar{\mathfrak{p}}\cdot
\mathfrak{a}$). Then
\[
\bigl(\varphi_{P}(\alpha_{1},\alpha_{2})\bigr)
=\bar{\mathfrak{p}}^{12},
\qquad
\bigl(\varphi_{\bar{P}}(\alpha_{1},\alpha_{2})\bigr)
=\mathfrak{p}^{12},
\]
and if $M\in H_{p}$, $M\notin GP$, $G\bar{P}$, then
$\varphi_{M}(\alpha_{1},\alpha_{2})$ is a unit.
\end{theorem}

\begin{proof}[Proof.]
Let $f$ be a positive integer such that $\mathfrak{p}^{f}$ is a principal
ideal in $K$. We have then $\mathfrak{p}^{f}=(\alpha)$, and
$\alpha\bar{\alpha}=N_{K/\mathbb{Q}}(\alpha)=p^{f}$, where $\alpha$ is an
integer of $K$. By the discussion preceding Theorem~\ref{thm:IV-1}, the
elements $P,\bar{P}$ are in $H_{p}$, and we can find
$P_{1}=P,P_{2},\ldots,P_{f}$ in $H_{p}$ such that
$P_{i}P_{i-1}\ldots P_{1}(\alpha_{1},\alpha_{2})$ is a basis of
$\mathfrak{p}^{i}\cdot\mathfrak{a}$ $(1\leq i\leq f)$ and, in particular,
$P_{f}\ldots P_{1}(\alpha_{1},\alpha_{2})=(\alpha\alpha_{1},\alpha\alpha_{2})$.
Put
\[
\lambda_{i}=\varphi_{P_{i}}\bigl(P_{i-1}\ldots P_{1}(\alpha_{1},\alpha_{2})\bigr)
=p^{12}\,\frac{\Delta(P_{i}\ldots P_{1}(\alpha_{1},\alpha_{2}))}
{\Delta(P_{i-1}\ldots P_{1}(\alpha_{1},\alpha_{2}))},\quad(1\leq i\leq f).
\]
Then,
\[
\prod_{i=1}^{f}\lambda_{i}
=p^{12f}\,\frac{\Delta(P_{f}\dots P_{1}(\alpha_{1},\alpha_{2}))}
{\Delta(\alpha_{1},\alpha_{2})}
=p^{12f}\,\frac{\Delta(\alpha\alpha_{1},\alpha\alpha_{2})}
{\Delta(\alpha_{1},\alpha_{2})}
=p^{12f}\alpha^{-12}
=\bar{\alpha}^{12}.
\]
As every $\lambda_{i}$ is an algebraic integer by
Lemma~\ref{lem:IV-2}, the ideal $(\lambda_{i})$ divides
$(\overline{\alpha}^{12})=\bar{\mathfrak{p}}^{12f}$. But, by the same
lemma, $(\lambda_{i})$ also divides $(p^{12})
=\mathfrak{p}^{12}\,\bar{\mathfrak{p}}^{12}$. Hence, $(\lambda_{i})$ is a
divisor of $\bar{\mathfrak{p}}^{12}$. Since
$\prod_{i=1}^{f}(\lambda_{i})=(\overline{\alpha}^{12})
=\bar{\mathfrak{p}}^{12f}$, it follows that $(\lambda_{i})
=\bar{\mathfrak{p}}^{12}$ for every $i$, and, in particular,
$\bigl(\varphi_{P}(\alpha_{1},\alpha_{2})\bigr)=(\lambda_{1})
=\bar{\mathfrak{p}}^{12}$. Similarly,
$\bigl(\varphi_{\bar{P}}(\alpha_{1},\alpha_{2})\bigr)
=\mathfrak{p}^{12}$. As $\mathfrak{p}\neq\bar{\mathfrak{p}}$, this shows
that $GP\neq G\bar{P}$, and it follows from
\[
\prod_{i=1}^{N}\bigl(\varphi_{i}(\alpha_{1},\alpha_{2})\bigr)=(p^{12}),
\]
that $\varphi_{M}(\alpha_{1},\alpha_{2})$ is a unit if $GM\neq GP$,
$G\bar{P}$.
\end{proof}

% ---------------------------------------------------------------
\section{A formal congruence}\label{sec:IV-formal}

Let $\xi$ and $\eta$ be power series (in particular polynomials) of
certain variables with integral algebraic coefficients and $\mathfrak{a}$
an ideal of integral numbers. As usual we write
\[
\xi\equiv\eta\mod\mathfrak{a}
\]
if every coefficient of $\xi-\eta$ is in $\mathfrak{a}$.

We consider $2(p+1)$ q-series $\xi_{i}(q)$, $\eta_{i}(q)$,
$(1\leq i\leq p+1)$, with the following properties:

(i) Their coefficients are integral numbers of $\mathbb{Q}(\zeta_{p})$;
$\xi_{p+1}$, $\eta_{p+1}$ have rational integral coefficients. The
$\xi_{i}$'s (resp.\ $\eta_{i}$'s) $(1\leq i\leq p)$ are permuted by the
automorphisms of the Galois group of $\mathbb{Q}(\zeta_{p})/\mathbb{Q}$,
and $\xi_{i}\equiv\xi_{j}$, $\eta_{i}\equiv\eta_{j}\bmod(1-\zeta_{p})$,
$(1\leq i,j\leq p)$.

(ii) They represent holomorphic functions in the upper half plane which
are permuted among themselves by $G$.

By \eqref{eq:IV-3} the $\varphi_{i}$'s satisfy these assumptions. The
same is true for the $j_{i}$'s because we have
\[
\left\{
\begin{aligned}
j_{i}(q)&=\zeta_{p}^{-i}q^{-1/p}\bigl(1+A(\zeta_{p}^{i}q^{1/p})\bigr)\\
j_{p+1}(q)&=q^{-p}\bigl(1+A(q^{p})\bigr)
\end{aligned}\right.
\tag{4}\label{eq:IV-4}
\]
where $A(q)$ has rational integral coefficients. In fact, these will be
the only functions satisfying (i), (ii) to be considered in the
applications.

With indeterminates $t,u$ we then put
\[
G_{p}(t,u,\xi_{i},\eta_{j})
=\sum_{j=1}^{j=p+1}\Bigl[(t-\xi_{j})
\prod_{\substack{k=1\\ k\neq j}}^{k=p+1}(u-\eta_{k})\Bigr].
\tag{5}\label{eq:IV-5}
\]

\begin{lemma}\label{lem:IV-4}
(a) $G_{p}(t,u,\xi_{i},\eta_{j})=\widetilde{G}_{p}(t,u,j)$ is a
polynomial in $t,u,j$ with rational integral coefficients.

(b) $G_{p}\bigl(t,u,\xi_{i}(q),\eta_{j}(q)\bigr)
\equiv(t-\xi_{p+1}(q))\bigl(u^{p}-\eta_{p}(q)^{p}\bigr)\bmod p$.
\end{lemma}

\begin{proof}[Proof.]
(a) is proved exactly as Theorem~1a of III. Since in
$\mathbb{Q}(\zeta_{p})$, the ideal $(p)$ is equal to
$(1-\zeta_{p})^{p-1}$ and since both sides of (b) have rational integral
coefficients, it is enough to show that both sides of (b) are congruent
mod $(1-\zeta_{p})$. By (i) above, we have $\xi_{i}=\xi_{j}$ and
$\eta_{i}=\eta_{j}\bmod(1-\zeta_{p})$ for $(1\leq j,i\leq p)$. Therefore
\[
G_{p}\equiv(t-\xi_{p+1})(u-\eta_{p})^{p}
+p\,(u-\eta_{p+1})(t-\xi_{p})(u-\eta_{p})^{p-1}\bmod(1-\zeta_{p}),
\]
and our contention follows from the fact that mod $p$, we have
$(u-\eta_{p})^{p}\equiv u^{p}-\eta_{p}^{p}$.
\end{proof}

% ---------------------------------------------------------------
\section{Applications}\label{sec:IV-applications}

The modular equation for the degree $p$ is $F_{p}(t,j)=0$ where
$F_{p}(t,j)=\prod_{i}(t-j_{i})$ is a polynomial in $t$ and $j$ with
rational integral coefficients (III, Theorem~1a). The following property
of $F_{p}$ is a classical result due to H.~Weber (Acta Mathematica 6,
1885, p.~390).

\begin{theorem}\label{thm:IV-2}
We have $F_{p}(t,j)\equiv(t-j^{p})(t^{p}-j)\bmod p$, the integer $p$
being prime.
\end{theorem}

\begin{proof}[Proof.]
From \eqref{eq:IV-4} it follows by Fermat's first theorem that
\[
\begin{aligned}
j_{p+1}(q)&\equiv j(q)^{p} &&\text{mod }p\\
j_{p}(q)^{p}&\equiv j(q) &&\text{mod }p
\end{aligned}
\tag{6}\label{eq:IV-6}
\]
In order to apply Lemma~\ref{lem:IV-4}, we put $t=u$,
$\xi_{i}=\eta_{i}=j_{i}$ $(1\leq i\leq p+1)$. Then
$\widetilde{G}_{p}=F_{p}(t,j)$, and, in view of \eqref{eq:IV-6}, the right
hand side in Lemma~\ref{lem:IV-4}b is congruent to
$(t-j^{p})(t^{p}-j)\bmod p$. Lemma~\ref{lem:IV-4} implies therefore that
the difference of the two sides in the formula of
Theorem~\ref{thm:IV-2}, when viewed as a power series in $t$ and $q$, has
all its coefficients in $p\mathbb{Z}$. By the q-expansion principle
(II §6), the same is then true when this difference is written as a
polynomial in $t$ and $j$.
\end{proof}

\begin{theorem}\label{thm:IV-3}
Let $K$ be an imaginary quadratic field, $\mathfrak{p}$ a prime ideal of
$K$ of degree $1$, unramified for $K/\mathbb{Q}$, and $\mathfrak{k}_{p}$
the ideal class containing $\mathfrak{p}$. Then for any ideal
$\mathfrak{k}$ of $K$, we have
\[
j(\mathfrak{k}_{p}^{-1}\cdot\mathfrak{k})
\equiv j(\mathfrak{k})^{N\mathfrak{p}}\mod\mathfrak{p}.
\tag{7}\label{eq:IV-7}
\]
\end{theorem}

\begin{proof}[Proof.]
Let $\mathfrak{a}$ be an ideal of $\mathfrak{k}$. Since
$\mathfrak{p}\cdot\bar{\mathfrak{p}}=(N\mathfrak{p})$, the ideal class
$\mathfrak{k}_{p}^{-1}$ contains $\bar{\mathfrak{p}}$ and our contention
is equivalent to
\[
j(\mathfrak{p}\cdot\mathfrak{a})\equiv j(\mathfrak{a})^{p}
\mod\mathfrak{p}\qquad(p=N\mathfrak{p}).
\tag{8}\label{eq:IV-8}
\]
Here $p$ is prime, and $\mathfrak{p}\neq\bar{\mathfrak{p}}$. Let
$\alpha_{1},\alpha_{2}$ be a base of $\mathfrak{a}$ and
$P,\bar{P}\in H_{p}$ be as in Theorem~\ref{thm:IV-1}. Furthermore, let
$c,d\leq p+1$ be the indices such that $P\in GM_{c}$,
$\bar{P}\in GM_{d}$. Then $c\neq d$ since $\mathfrak{p}\neq\bar{\mathfrak{p}}$,
and
\[
j(\bar{\mathfrak{p}}\cdot\mathfrak{a})=j_{\bar{p}}(\mathfrak{a})
=j_{d}(\mathfrak{a})=j_{d}(\alpha_{1}/\alpha_{2}).
\tag{9}\label{eq:IV-9}
\]

We now apply Lemma~\ref{lem:IV-4} to the case where $\xi_{i}=j_{i}$,
$\eta_{i}=\varphi_{i}$ $(1\leq i\leq p+1)$ and $t=j^{p}$. This gives
\[
G_{p}\bigl(j(q)^{p},u,j_{i}(q),\varphi_{k}(q)\bigr)
\equiv(j^{p}-j_{p+1})(u^{p}-\varphi_{p}^{p})\bmod p,
\]
which, together with \eqref{eq:IV-6}, shows that
\[
G_{p}\bigl(j(q)^{p},u,j_{i}(q),\varphi_{k}(q)\bigr)\equiv 0\bmod p,
\]
$G_{p}$ being considered as a power series in $u$ and $q$. By the
q-expansion principle (II §6), we have then
\[
\widetilde{G}_{p}(j^{p},u,j)\equiv 0\bmod p,
\tag{10}\label{eq:IV-10}
\]
where (cf.\ Lemma~\ref{lem:IV-4}), $\widetilde{G}_{p}$ is $G_{p}$
considered as a polynomial in $j$ and $u$.

Let now $q_{0}=\exp(2\pi i\alpha_{1}/\alpha_{2})$. The numbers
$j(q_{0}),j_{i}(q_{0}),\varphi_{i}(q_{0})$ are algebraic integers. Since
$G_{p}(j^{p},u,j_{i},\varphi_{j})$ may be written as a polynomial in $t$
and $j$ with coefficients in $p\mathbb{Z}$, its value for $q=q_{0}$ is a
polynomial in $u$, whose coefficients are algebraic integers divisible by
$p$. Thus
\[
G_{p}\bigl(j(\mathfrak{a})^{p},u,j_{i}(\mathfrak{a}),
\varphi_{k}(\alpha_{1},\alpha_{2})\bigr)\equiv 0\bmod p.
\tag{11}\label{eq:IV-11}
\]

Let us now put $u=\varphi_{\bar{P}}(\alpha_{1},\alpha_{2})
=\varphi_{d}(\alpha_{1},\alpha_{2})$. Then by \eqref{eq:IV-11} and the
definition of $G_{p}$, we get
\[
\bigl(j(\mathfrak{a})^{p}-j_{d}(\mathfrak{a})\bigr)
\prod_{i\neq d}\bigl(\varphi_{d}(\alpha_{1},\alpha_{2})
-\varphi_{i}(\alpha_{1},\alpha_{2})\bigr)\equiv 0\bmod p.
\tag{12}\label{eq:IV-12}
\]

Since $\mathfrak{p}$ divides $(p)$, we have a fortiori
\[
\bigl(j(\mathfrak{a})^{p}-j_{d}(\mathfrak{a})\bigr)
\prod_{i\neq d}\bigl(\varphi_{d}(a_{1},a_{2})
-\varphi_{i}(a_{1},a_{2})\bigr)\equiv 0\bmod\mathfrak{p}.
\tag{13}\label{eq:IV-13}
\]
By Theorem~\ref{thm:IV-1},
\[
\prod_{i\neq d}\bigl(\varphi_{d}(\alpha_{1},\alpha_{2})
-\varphi_{i}(\alpha_{1},\alpha_{2})\bigr)
\equiv(-1)^{p}\prod_{i\neq d}\varphi_{i}(\alpha_{1},\alpha_{2})
\bmod\mathfrak{p}
\]
and the right hand side generates the ideal $\bar{\mathfrak{p}}^{12}$,
which is prime to $\mathfrak{p}$. Therefore, \eqref{eq:IV-13} implies
that
\[
j(\mathfrak{a})^{p}-j_{d}(\mathfrak{a})\equiv 0\bmod\mathfrak{p}
\]
and this, in view of \eqref{eq:IV-9}, is the congruence
\eqref{eq:IV-8}.
\end{proof}

\subsection*{Remarks.}

\noindent
(1) We have $j_{p+1}(z)=j(p.z)$, so that the first congruence in
\eqref{eq:IV-6} is a congruence between the q-developments of $j(z)$ and
$j(pz)$, that is between $j(q)$ and $j(q^{p})$. It may be viewed as a
formal analogue to Theorem~\ref{thm:IV-3}. In fact, for a suitable base
$\alpha_{1},\alpha_{2}$ of $\mathfrak{a}$ the numbers
$p\alpha_{1},\alpha_{2}$ form a base of $\bar{\mathfrak{p}}\cdot
\mathfrak{a}$ so that \eqref{eq:IV-8} may be written
\[
j(\mathfrak{a})^{p}\equiv j_{p+1}(\mathfrak{a})\bmod\mathfrak{p}.
\tag{14}\label{eq:IV-14}
\]
\begin{sloppypar}
However, one cannot of course simply derive \eqref{eq:IV-14} out of
\eqref{eq:IV-6} by putting $q=q_{0}=\exp(2\pi i\alpha_{1}/\alpha_{2})$,
because both sides are infinite series in $q$ (and as a matter of fact,
this would lead to a congruence mod $\bar{\mathfrak{p}}$, not only mod
$\mathfrak{p}$, which is not true in general). However, a formal
congruence between q-series with coefficients in an additive group $H$ of
algebraic integers leads to a correct congruence for $q=q_{0}$ if both
sides represent integral modular functions, because they can then be
considered as polynomials in $j$ with coefficients in $H$ by the
q-expansion principle, and $j(q_{0})$ is an algebraic integer. It is this
fact which has allowed us to derive \eqref{eq:IV-11} from
Lemma~\ref{lem:IV-4}.
\end{sloppypar}

\medskip\noindent
(2) M.~Eichler, Math.\ Zeitschrift 64 (1956), 229--242
\cite{eichler1956quadratische} has given a proof of the main theorems of
complex multiplication which uses only $j$, and not $\Delta$. However, he
does not obtain Theorem~\ref{thm:IV-3} in full, but shows instead that one
has either Theorem~\ref{thm:IV-3} or
$j(\mathfrak{k}_{p}\cdot\mathfrak{k})\equiv j(\mathfrak{k})^{p}\bmod
\mathfrak{p}$.
